\section{KẾT LUẬN VÀ KIẾN NGHỊ}

\subsection{Kết luận}
Qua quá trình tính toán thiết kế kết cấu áo đường mềm cho tuyến đường cấp II đồng bằng với các số liệu đầu vào đã cho, bài làm rút ra các kết luận sau:

\begin{enumerate}
    \item \textbf{Lưu lượng xe:} Tuyến đường có lưu lượng xe quy đổi tích lũy $N_e \approx 3.86 \times 10^6$ trục xe tiêu chuẩn 100 kN trong 15 năm. Đây là tuyến đường có lưu lượng lớn.
    \item \textbf{Kết cấu áo đường:} Để đảm bảo cường độ chung $E_{yc} \ge 210$ MPa và các điều kiện chịu cắt, chịu kéo uốn, kết cấu áo đường được đề xuất gồm 4 lớp vật liệu trên nền đất đầm chặt ($K \ge 0.98$):
    \begin{itemize}
        \item Lớp mặt: 12 cm Bê tông nhựa (5cm BTN C12.5 + 7cm BTN C19).
        \item Lớp móng: 50 cm Cấp phối đá dăm (20cm CPĐD I + 30cm CPĐD II).
    \end{itemize}
    \item Kết cấu này đảm bảo khả năng khai thác ổn định trong suốt thời hạn thiết kế 15 năm với tốc độ tăng trưởng xe 6\%/năm.
\end{enumerate}

\subsection{Kiến nghị}
\begin{itemize}
    \item Cần chú ý công tác thi công và nghiệm thu từng lớp kết cấu, đặc biệt là độ chặt của các lớp móng cấp phối đá dăm.
    \item Nên xem xét bố trí hệ thống thoát nước dọc và ngang hợp lý để bảo vệ nền đường, duy trì cường độ nền đất $E_0$.
\end{itemize}
